\documentclass[conference]{IEEEtran}

\usepackage{amsmath,amssymb}
\usepackage{booktabs}
\usepackage{multirow}
\usepackage{graphicx}
\usepackage{xcolor}
\usepackage{url}

\newcommand{\blank}{\varnothing}
\newcommand{\eoc}{\mathrm{EOC}}
\newcommand{\logsumexp}{\operatorname{logsumexp}}

\title{Chunked Aligners: Alignment-Free Streaming Speech Recognition with Bounded Per-Chunk Emissions}

\author{
  \IEEEauthorblockN{Anonymous Authors}
  \IEEEauthorblockA{Anonymous Affiliation\\
  \texttt{anonymous@example.com}}
}

\begin{document}
\maketitle

\begin{abstract}
Streaming automatic speech recognition requires models that can emit tokens with
bounded latency while remaining simple and memory efficient to train. Prior
chunkwise aligner models impose a local monotonic alignment structure, but rely
on external timestamp supervision or alignments produced by another model. We
propose the \emph{Chunked Aligner}, a full-sum training objective that learns
chunk-level alignments directly from transcripts. The encoder sequence is
partitioned into fixed-size chunks. Within each chunk, a token emission advances
one local frame and a blank symbol acts as an end-of-chunk transition to the next
chunk. A forward-backward algorithm sums over all possible transcript
segmentations across chunks, analogous to RNN-T and CTC, without requiring
timestamps. We implement both autoregressive and non-autoregressive variants.
The autoregressive model reuses an RNN-T-style prediction network and joint,
whereas the non-autoregressive model removes the predictor and joint entirely and
uses a per-frame projection head, reducing the dominant activation tensor from
$O(BTUV)$ to $O(BTV)$. We further study two methods for controlling the maximum
number of tokens emitted per acoustic chunk: a learned token extraction layer and
a parameter-free first-$k$ frame selection. CUDA kernels are implemented with
Numba and validated against PyTorch reference losses and brute-force path
enumeration. The resulting model provides a simple streaming transducer
objective that preserves exact marginalization over alignments while giving
explicit control over chunk-local emission capacity.
\end{abstract}

\begin{IEEEkeywords}
streaming ASR, transducer, aligner, forward-backward, chunked decoding,
non-autoregressive ASR
\end{IEEEkeywords}

\section{Introduction}

End-to-end automatic speech recognition (ASR) systems must often operate under
latency constraints. RNN-Transducer (RNN-T) \cite{graves2012rnnt} and CTC
\cite{graves2006ctc} are popular because they marginalize over monotonic
alignments instead of requiring frame-level labels. However, RNN-T training
materializes a joint distribution over acoustic time and output label position,
and streaming behavior is usually controlled indirectly through encoder
lookahead, decoding constraints, or duration modeling.

Aligner-Encoder models \cite{stooke2025aligner} suggest a different view: make
the encoder self-transducing by pairing encoder frames and output tokens more
directly. Chunkwise aligner approaches further constrain emissions to local
chunks for streaming recognition, but existing formulations commonly depend on
timestamps or forced alignments from a teacher model. Such external supervision
can be inconvenient, brittle, and suboptimal when the desired alignment behavior
differs from the teacher.

This paper describes a chunked aligner objective that learns the chunk
segmentation directly from transcript supervision. Each encoder chunk may emit a
bounded number of tokens. A blank symbol serves as an end-of-chunk (EOC) signal:
emitting a blank skips the remainder of the current chunk and advances to the
next chunk. Non-blank emissions advance one frame within the chunk. The training
criterion sums the probabilities of all paths that produce the target transcript,
using a dynamic program similar in spirit to CTC and RNN-T.

The main contributions are:
\begin{itemize}
  \item We define a chunk-local transducer lattice that learns alignments without
  external timestamps while enforcing a bounded number of emissions per chunk.
  \item We derive an exact forward-backward loss for both autoregressive (AR) and
  non-autoregressive (NAR) chunked aligner models.
  \item We introduce two ways to control per-chunk capacity: learned token
  extraction by cross-attention and a parameter-free first-$k$ frame selection.
  \item We implement CUDA/Numba kernels for the objective and validate them
  against PyTorch autograd references, brute-force enumeration, and a
  ``trivial-join'' equivalence between NAR and AR models.
\end{itemize}

\section{Chunked Aligner Lattice}

Let $x$ denote an input utterance and let the encoder produce acoustic
representations
\begin{equation}
  H = (h_0, h_1, \ldots, h_{T-1}), \quad h_t \in \mathbb{R}^d.
\end{equation}
The target transcript is
\begin{equation}
  y = (y_0, y_1, \ldots, y_{U-1}), \quad y_u \in \mathcal{V},
\end{equation}
where $\mathcal{V}$ is the vocabulary. A special blank token $\blank$ is added to
represent an end-of-chunk transition.

The encoder timeline is partitioned into chunks of $C$ acoustic frames. For a
time index $t$, define the chunk index and local offset as
\begin{equation}
  n(t) = \left\lfloor \frac{t}{C} \right\rfloor, \qquad
  j(t) = t \bmod C.
\end{equation}
Within a chunk, non-blank emissions advance diagonally:
\begin{equation}
  (t,u) \rightarrow (t+1,u+1), \quad \text{emit } y_u.
\end{equation}
Blank emissions act as EOC transitions:
\begin{equation}
  (t,u) \rightarrow ((n(t)+1)C, u), \quad \text{emit } \blank.
\end{equation}
Thus each chunk can emit at most $C$ tokens in the basic formulation. Later we
replace $C$ by a smaller effective value using token extraction or first-$k$
selection.

\section{Autoregressive Model}

The AR model uses an RNN-T-style prediction network. Given the previous emitted
tokens $y_{<u}$, the prediction network produces a state $g_u$. A joint network
combines acoustic and prediction states:
\begin{equation}
  z_{t,u} = J(h_t, g_u),
  \qquad
  p_{t,u}(a) = \mathrm{softmax}(z_{t,u})_a,
\end{equation}
for $a \in \mathcal{V} \cup \{\blank\}$. The implementation passes raw logits to
the loss wrapper, applies a single log-softmax in PyTorch, and uses CUDA kernels
over normalized log-probabilities.

\subsection{Forward Recursion}

Let $\alpha(t,u)$ be the log probability of arriving at encoder frame $t$ after
having emitted exactly $u$ labels. The base condition is
\begin{equation}
  \alpha(0,0) = 0,
\end{equation}
and all other states are initialized to $-\infty$. For $t>0$ or $u>0$, the state
receives two types of contributions. The token contribution is
\begin{equation}
  \alpha(t-1,u-1) + \log p_{t-1,u-1}(y_{u-1}),
\end{equation}
valid when $t \ge 1$ and $u \ge 1$. At a chunk boundary, $j(t)=0$ and
$n(t)\ge1$, blank transitions from any frame in the previous chunk can land at
$(t,u)$:
\begin{equation}
  \alpha(t',u) + \log p_{t',u}(\blank),
  \quad t' \in [(n(t)-1)C, n(t)C-1].
\end{equation}
Therefore
\begin{align}
  \alpha(t,u) =
  \logsumexp \Big(&
  \{\alpha(t-1,u-1)+\log p_{t-1,u-1}(y_{u-1})\} \nonumber\\
  &\cup
  \{\alpha(t',u)+\log p_{t',u}(\blank)\}_{t'} \Big),
\end{align}
where invalid terms are omitted.

\subsection{Termination}

Let $N=\lceil T/C\rceil$ and let $s=(N-1)C$ be the start index of the last chunk.
An alignment is complete if either all labels have been emitted and a blank is
emitted somewhere in the final chunk, or the last label is emitted on the final
encoder frame. The total log likelihood is
\begin{align}
  \log P(y|x) = \logsumexp \Big(
  &\{\alpha(t,U) + \log p_{t,U}(\blank) : s \le t < T\} \nonumber\\
  &\cup \{\alpha(T-1,U-1) + \log p_{T-1,U-1}(y_{U-1})\}
  \Big).
\end{align}
The loss is the negative log likelihood.

\subsection{Backward Recursion and Gradients}

The backward variable $\beta(t,u)$ is defined as the log probability of
completing the transcript from state $(t,u)$. It follows the reverse of the same
two arc types: a token arc to $(t+1,u+1)$ and a blank arc to the next chunk start.
For normalized log-probabilities, the gradient of $-\log P(y|x)$ with respect to
an arc log-probability is the negative posterior probability of that arc. The
Numba kernel computes these posterior counts from $\alpha$, $\beta$, and the
forward log likelihood, then PyTorch autograd applies the log-softmax Jacobian to
obtain gradients with respect to raw logits.

\section{Non-Autoregressive Chunked Aligner}

The NAR model removes the prediction network and joint network. Each acoustic
representation is mapped directly to vocabulary logits:
\begin{equation}
  z_t = W h_t + b, \qquad
  p_t(a) = \mathrm{softmax}(z_t)_a.
\end{equation}
The same chunked lattice is used, but arc scores no longer depend on $u$:
\begin{align}
  \text{token arc:}\quad & \log p_{t}(y_u),\\
  \text{blank arc:}\quad & \log p_{t}(\blank).
\end{align}
The forward and backward dynamic programs remain over $(t,u)$ states, but the
large activation tensor changes from $[B,T,U,V]$ to $[B,T,V]$. This removes the
dominant joint activation and the $T \times U$ joint computation. The gradient
with respect to a frame-level log-probability is the sum over all predictor
positions $u$ of the corresponding AR arc posteriors. In implementation we
validate this relationship by constructing a ``trivial join'' tensor that
broadcasts NAR logits over the $u$ axis and checking equality of both loss and
gradient with the AR objective.

\section{Controlling Per-Chunk Capacity}

The basic chunked aligner permits at most $C$ emissions per chunk. For streaming
recognition it can be useful to decouple the acoustic context duration from the
maximum number of emitted tokens. We implemented two mechanisms.

\subsection{Learned Token Extraction}

For an acoustic chunk
\begin{equation}
  H_n = (h_{nC}, \ldots, h_{nC+C-1}),
\end{equation}
we introduce $K$ trainable query vectors
\begin{equation}
  Q = (q_1, \ldots, q_K), \quad q_i \in \mathbb{R}^d.
\end{equation}
The queries cross-attend to the $C$ encoder frames in the chunk:
\begin{equation}
  \tilde{H}_n = \mathrm{CrossAttn}(Q, H_n, H_n),
  \qquad \tilde{H}_n \in \mathbb{R}^{K \times d}.
\end{equation}
Padding masks prevent attention to invalid frames in the final partial chunk.
The chunked aligner lattice then runs over the extracted sequence
$\tilde{H}$, with effective chunk size $K$. This learned bottleneck can reduce
joint memory and explicitly set the maximum number of tokens per acoustic chunk.

\subsection{First-$k$ Frame Selection}

As a parameter-free alternative, we keep only the first $k$ frames from each
acoustic chunk:
\begin{equation}
  \tilde{H}_n = (h_{nC}, h_{nC+1}, \ldots, h_{nC+k-1}),
  \qquad k \le C.
\end{equation}
For a partial final chunk of length $r$, the number of retained frames is
$\min(k,r)$. This gives the same emission-capacity control as token extraction
without adding parameters or attention computation. In our implementation this
option is controlled by \texttt{first\_k\_frames\_per\_chunk}; the default value
$-1$ disables it. It is mutually exclusive with learned token extraction.

\section{Implementation}

The loss is implemented in three forms.
\begin{enumerate}
  \item A loop-based PyTorch reference implementation used for clarity and
  gradient checks.
  \item A brute-force path enumerator used only for tiny examples.
  \item CUDA kernels written in Numba for forward variables, backward variables,
  and sparse arc-posterior gradients.
\end{enumerate}
The AR kernel operates on log-probabilities with shape $[B,T,U,V]$. The NAR
kernel operates on $[B,T,V]$ and accumulates gradient contributions across $u$.
Both kernels receive already normalized log-probabilities; the log-softmax is
kept in PyTorch to simplify the CUDA implementation while preserving exact
gradients through autograd.

The method is integrated into the NeMo RNN-T model lifecycle. Setting
\texttt{model.loss\_type=chunked\_aligner} activates the chunked-aligner loss and
greedy decoding. Setting \texttt{model.chunked\_aligner.nar=true} constructs the
NAR architecture and avoids instantiating unused prediction and joint modules.
Token extraction and first-$k$ selection are applied immediately after the
encoder, before the AR joint or NAR projection head.

\section{Inference Algorithms}

Greedy decoding replays the chunked lattice. In the AR model, at each visited
frame the joint network predicts a token or blank. A blank skips the remainder of
the current chunk; a token is emitted, fed into the predictor, and the local frame
index advances by one. In the NAR model, all per-frame logits are computed in one
batched projection, and decoding reduces to an argmax scan with the same blank
chunk-jump rule.

The current AR greedy decoder is a per-utterance reference implementation. It is
correct but not optimized for throughput. The NAR decoder is naturally batched in
the projection step because it has no predictor state.

\subsection{Common Encoder and Chunk Reduction}

All chunked variants first encode the waveform and optionally reduce each
acoustic chunk. The reduction step determines the effective number of lattice
frames per chunk. With no reduction, this value is the acoustic chunk size $C$.
With token extraction or first-$k$ selection, it becomes $K$ or $k$ respectively.
Algorithm~\ref{alg:reduce} summarizes this shared preprocessing.

\begin{figure}[t]
\small
\begin{tabbing}
\hspace{0.03in}\=\hspace{0.20in}\=\hspace{0.20in}\=\kill
\textbf{Algorithm 1: Common encoder and chunk reduction}\\
\textbf{Input:} waveform $x$, acoustic chunk size $C$\\
\textbf{Output:} reduced sequence $\tilde{H}$, lengths $\tilde{T}$, effective chunk size $\tilde{C}$\\
$H,T \leftarrow \mathrm{Encode}(x)$\\
\textbf{if} \texttt{token\_extraction\_size} $=K \ge 1$ \textbf{then}\\
\> $\tilde{H} \leftarrow [\,\mathrm{CrossAttn}(Q,H_n,H_n)\,]_{n=0}^{\lceil T/C\rceil-1}$\\
\> $\tilde{T} \leftarrow \sum_n \min(K, K\cdot \mathbb{1}[|H_n|>0])$ \\
\> $\tilde{C} \leftarrow K$\\
\textbf{else if} \texttt{first\_k\_frames\_per\_chunk} $=k \ge 1$ \textbf{then}\\
\> $\tilde{H} \leftarrow [\,h_{nC},\ldots,h_{nC+\min(k,|H_n|)-1}\,]_n$\\
\> $\tilde{T} \leftarrow \sum_n \min(k, |H_n|)$\\
\> $\tilde{C} \leftarrow k$\\
\textbf{else}\\
\> $\tilde{H} \leftarrow H,\quad \tilde{T} \leftarrow T,\quad \tilde{C} \leftarrow C$\\
\textbf{return} $\tilde{H},\tilde{T},\tilde{C}$
\end{tabbing}
\caption{Shared inference preprocessing for chunked aligner variants. $H_n$
denotes the $n$-th acoustic chunk.}
\label{alg:reduce}
\end{figure}

\subsection{Autoregressive Chunked Aligner}

The AR chunked decoder keeps an RNN-T-style prediction state. It visits at most
one frame per local emission step. When a blank is emitted, the decoder jumps to
the next chunk boundary without changing the prediction state. When a token is
emitted, the token is appended to the hypothesis and fed to the predictor.

\begin{figure}[t]
\small
\begin{tabbing}
\hspace{0.03in}\=\hspace{0.20in}\=\hspace{0.20in}\=\kill
\textbf{Algorithm 2: Chunked AR greedy inference}\\
\textbf{Input:} waveform $x$\\
\textbf{Output:} token sequence $\hat{y}$\\
$H,T,C \leftarrow \mathrm{ReduceChunks}(\mathrm{Encode}(x))$\\
$\hat{y} \leftarrow [\,]$\\
$g,s \leftarrow \mathrm{Predictor}(\mathrm{SOS}, \varnothing)$\\
$t \leftarrow 0$\\
\textbf{while} $t < T$ \textbf{do}\\
\> $z \leftarrow \mathrm{Joint}(H_t, g)$\\
\> $a \leftarrow \arg\max_{v \in \mathcal{V}\cup\{\blank\}} z_v$\\
\> \textbf{if} $a=\blank$ \textbf{then}\\
\>\> $t \leftarrow (\lfloor t/C \rfloor + 1)C$ \hfill // EOC jump\\
\> \textbf{else}\\
\>\> $\hat{y}.\mathrm{append}(a)$\\
\>\> $g,s \leftarrow \mathrm{Predictor}(a,s)$\\
\>\> $t \leftarrow t+1$\\
\textbf{return} $\hat{y}$
\end{tabbing}
\caption{Greedy inference for the autoregressive chunked aligner.}
\label{alg:chunked-ar}
\end{figure}

\subsection{Non-Autoregressive Chunked Aligner}

The NAR decoder has no predictor state. It first computes all frame-level logits
with a single projection head and then performs the same chunk-jump scan over the
argmax labels. This gives the same streaming control flow as the AR model while
avoiding per-label prediction-network calls.

\begin{figure}[t]
\small
\begin{tabbing}
\hspace{0.03in}\=\hspace{0.20in}\=\hspace{0.20in}\=\kill
\textbf{Algorithm 3: Chunked NAR greedy inference}\\
\textbf{Input:} waveform $x$\\
\textbf{Output:} token sequence $\hat{y}$\\
$H,T,C \leftarrow \mathrm{ReduceChunks}(\mathrm{Encode}(x))$\\
$Z \leftarrow \mathrm{Linear}(H)$ \hfill // $Z \in \mathbb{R}^{T \times |\mathcal{V}\cup\{\blank\}|}$\\
$A_t \leftarrow \arg\max_v Z_{t,v}$ for all $t$\\
$\hat{y} \leftarrow [\,]$\\
$t \leftarrow 0$\\
\textbf{while} $t < T$ \textbf{do}\\
\> $a \leftarrow A_t$\\
\> \textbf{if} $a=\blank$ \textbf{then}\\
\>\> $t \leftarrow (\lfloor t/C \rfloor + 1)C$ \hfill // EOC jump\\
\> \textbf{else}\\
\>\> $\hat{y}.\mathrm{append}(a)$\\
\>\> $t \leftarrow t+1$\\
\textbf{return} $\hat{y}$
\end{tabbing}
\caption{Greedy inference for the non-autoregressive chunked aligner.}
\label{alg:chunked-nar}
\end{figure}

\subsection{Nonchunked AR Aligner Baseline}

For comparison, the original AR aligner performs a left-to-right one-to-one
walk over encoder frames. There is no blank or chunk jump; decoding stops when an
EOS token is predicted or the encoder sequence ends.

\begin{figure}[t]
\small
\begin{tabbing}
\hspace{0.03in}\=\hspace{0.20in}\=\hspace{0.20in}\=\kill
\textbf{Algorithm 4: Nonchunked AR aligner greedy inference}\\
\textbf{Input:} waveform $x$\\
\textbf{Output:} token sequence $\hat{y}$\\
$H,T \leftarrow \mathrm{Encode}(x)$\\
$\hat{y} \leftarrow [\,]$\\
$g,s \leftarrow \mathrm{Predictor}(\mathrm{SOS}, \varnothing)$\\
\textbf{for} $t=0,\ldots,T-1$ \textbf{do}\\
\> $z \leftarrow \mathrm{AlignerJoint}(H_t,g)$\\
\> $a \leftarrow \arg\max_{v \in \mathcal{V}\cup\{\mathrm{EOS}\}} z_v$\\
\> \textbf{if} $a=\mathrm{EOS}$ \textbf{then break}\\
\> $\hat{y}.\mathrm{append}(a)$\\
\> $g,s \leftarrow \mathrm{Predictor}(a,s)$\\
\textbf{return} $\hat{y}$
\end{tabbing}
\caption{Greedy inference for the nonchunked AR aligner baseline.}
\label{alg:aligner-ar}
\end{figure}

\section{Experimental Plan}

We evaluate on English and Mandarin ASR settings. For Mandarin AISHELL-2, the
models use a character vocabulary and the same encoder configuration across
objectives. We compare:
\begin{itemize}
  \item a TDT transducer baseline;
  \item the nonchunked AR Aligner-Encoder;
  \item chunked AR aligners with and without token extraction;
  \item chunked NAR aligners with and without token extraction;
  \item chunked AR aligners using first-$k$ frame selection.
\end{itemize}
All models are trained from scratch under matched data, optimizer, and schedule
settings. Metrics are character error rate (CER) for Mandarin and word error
rate (WER) for English. We additionally report decoding speed and memory usage,
since the NAR and first-$k$ variants are designed to reduce training and
inference cost.

\section{Discussion}

The chunked aligner provides explicit control over streaming latency and local
emission capacity while preserving exact marginalization over alignments. Unlike
timestamp-supervised chunkwise aligners, it does not require an external aligner
or teacher model. The AR variant retains the expressiveness of a prediction
network. The NAR variant trades autoregressive dependence for a substantial
memory reduction, making it attractive when vocabulary size or transcript length
makes the joint tensor expensive.

The two capacity-control mechanisms have complementary tradeoffs. Learned token
extraction allows each acoustic chunk to be summarized by trainable queries, but
adds parameters and an attention module. First-$k$ selection is parameter-free and
directly preserves early chunk frames, but may discard useful information near the
end of the chunk. Empirically comparing these two mechanisms can clarify whether
the learned bottleneck is necessary or whether most of the benefit comes from the
capacity constraint itself.

\section{Conclusion}

We presented the Chunked Aligner, a streaming ASR objective that learns
chunk-level alignments from transcripts by summing over all valid paths in a
chunk-constrained lattice. The method generalizes aligner-style self-transduction
to streaming settings without timestamp supervision. We implemented exact
forward-backward losses for AR and NAR variants, together with practical
capacity-control mechanisms based on learned token extraction and first-$k$ frame
selection. The NAR formulation removes the joint tensor and provides a simple,
memory-efficient baseline for future streaming transducer research.

\begin{thebibliography}{9}

\bibitem{graves2006ctc}
A. Graves, S. Fernandez, F. Gomez, and J. Schmidhuber,
``Connectionist temporal classification: Labelling unsegmented sequence data with
recurrent neural networks,''
in \emph{Proc. ICML}, 2006.

\bibitem{graves2012rnnt}
A. Graves,
``Sequence transduction with recurrent neural networks,''
in \emph{ICML Workshop on Representation Learning}, 2012.

\bibitem{gulati2020conformer}
A. Gulati, J. Qin, C.-C. Chiu, N. Parmar, Y. Zhang, J. Yu, W. Han,
S. Wang, Z. Zhang, Y. Wu, and R. Pang,
``Conformer: Convolution-augmented Transformer for speech recognition,''
in \emph{Proc. Interspeech}, 2020.

\bibitem{stooke2025aligner}
A. Stooke, et al.,
``Aligner-Encoders: Self-Attention Transformers Can Be Self-Transducers,''
2025.

\bibitem{wenet2022}
B. Zhang, D. Wu, C. Yang, X. Chen, Z. Peng, X. Wang, Z. Yao, X. Wang,
F. Yu, L. Xie, and X. Lei,
``WeNet 2.0: More productive end-to-end speech recognition toolkit,''
in \emph{Proc. Interspeech}, 2022.

\bibitem{aishell2}
J. Du, X. Na, X. Liu, and H. Bu,
``AISHELL-2: Transforming Mandarin ASR research into industrial scale,''
arXiv:1808.10583, 2018.

\end{thebibliography}

\end{document}
